\documentclass[aps,pra,reprint,amsmath,amssymb,longbibliography,nofootinbib]{revtex4-2}
\usepackage[T1]{fontenc}
\usepackage{bm}
\usepackage{mathtools}
\usepackage{hyperref}
\hypersetup{colorlinks=true,citecolor=blue,linkcolor=blue,urlcolor=blue,
  pdftitle={Three-copy universal qubit purification under complete stabilizer preservation},
  pdfauthor={Xiaojian Zeng}}
\newcommand{\Tr}{\operatorname{Tr}}
\newcommand{\id}{\operatorname{id}}
\newcommand{\cE}{\mathcal E}
\newcommand{\cB}{\mathcal B}
\newcommand{\cA}{\mathcal A}
\newcommand{\cC}{\mathcal C}
\newcommand{\Qh}{\widehat Q}
\newcommand{\Rh}{\widehat R}
\newcommand{\Cs}{\mathsf C}
\newtheorem{theorem}{Theorem}
\newtheorem{lemma}{Lemma}
\begin{document}
\title{Three-copy universal qubit purification under complete stabilizer preservation}
\author{Xiaojian Zeng}
\email{xiaojianzeng@foxmail.com}
\affiliation{PBC School of Finance, Tsinghua University, Beijing 100190, China}
\date{September 12, 2026}
\begin{abstract}
We determine the optimal universal purification fidelity of three independently depolarized copies of an unknown pure qubit when the successful operation preserves the stabilizer positive cone completely. For every noise strength $0<\delta<1$ and every prescribed average success probability $0<s\leq1$, the success-weighted conditional Haar fidelity is exactly $1-\delta/2$. Keeping one input and accepting independently attains the bound. The proof expresses the fidelity witness as an exact cubic interpolation between a stabilizer Pauli-correlation inequality and a negative rank-five projector. It gives an analytic certificate for the three-copy case supported numerically in earlier work. Separate exact constructions of all four-qubit stabilizer projectors corroborate the certificate. The result concerns the specified successful branches and Haar objective, without assuming state-independent success or Clifford-circuit realizability of every admissible map.
\end{abstract}
\maketitle

\section{Introduction}

Universal quantum state purification asks whether several noisy copies of the same unknown target can yield an output of higher fidelity. A resource restriction on the operation is distinct from a restriction on the target state: the unknown qubit need not itself be a stabilizer state. This distinction matters when purification is intended to work over the full pure-state ensemble.

He \emph{et al.}\ proved a two-copy no-go theorem for qubits under completely stabilizer-preserving operations and reported numerical evidence for three and four copies~\cite{He2025NoGo}. In the fixed preprint version of that work, the analytic qubit result is Theorem~2, whereas the higher-copy evidence appears in Table~I. The subsequent two-copy resource laws of He, Xiong, and Wang quantify the nonstabilizerness needed for a fidelity gain~\cite{He2026Resource}. Their qubit robustness law and its no-go corollary remain two-to-one statements. Neither a two-copy upper bound nor the option of discarding an additional input gives a three-copy upper bound.

Here we give a direct analytic certificate for three inputs. With the Bloch-vector shrinkage $t=1-\delta$, the decisive operator identity is
\begin{equation}
 H(t)=\frac{t(1-t^2)}{48}(\cB-3I)-\frac{t^3}{5}P.
 \label{eq:certificate-intro}
\end{equation}
The nine-term Pauli operator $\cB$ has expectation at most three on every pure four-qubit stabilizer state, while $P$ is a rank-five positive projector. Equation~\eqref{eq:certificate-intro} therefore bounds every allowed successful branch throughout the noise interval. Its use requires the actual ensemble-averaged success probability; the trace of a normalized Choi operator is not a replacement for that probability.

Our literature comparison uses version~2 of Ref.~\cite{He2025NoGo} and version~1 of Ref.~\cite{He2026Resource} on arXiv. The published bibliographic record of the first work is given in Ref.~\cite{He2025NoGo}, but its final journal body and supplement were not available for comparison. We therefore make no absolute priority claim. The contribution stated here is the explicit three-copy analytic certificate and the resulting full-parameter operational bound.

\section{Model, objective, and successful-branch cone}
\label{sec:model}

Let $\psi=|\psi\rangle\langle\psi|$ be distributed according to normalized Haar measure on pure qubits. Each input undergoes the same local depolarizing channel,
\begin{align}
 \mathcal D_\delta(X)&=(1-\delta)X+\frac{\delta}{2}\Tr(X)I_2,
 \label{eq:noise}\\
 \rho_t(\psi)&=\mathcal D_\delta(\psi)
             =\frac{I_2+t\bm r\cdot\bm\sigma}{2},
 \qquad t=1-\delta,
 \label{eq:rho}
\end{align}
where $|\bm r|=1$. The joint input is $\rho_t(\psi)^{\otimes3}$: the target is shared and the noise acts independently on the three copies.

Write $\Cs_n$ for the nonnegative cone generated by pure $n$-qubit stabilizer projectors. A successful branch is a completely positive, trace-nonincreasing map $\cE:M_8\to M_2$ such that
\begin{equation}
  (\id_k\otimes\cE)(\Cs_{k+3})\subseteq\Cs_{k+1}
  \quad\text{for every integer }k\geq0,
  \label{eq:complete}
\end{equation}
where $\id_k$ acts on $k$ auxiliary qubits. The cone includes subnormalized states and the zero operator. This definition imposes complete preservation, including arbitrary finite auxiliary systems. It does not identify the class with a set of implementable Clifford circuits. The distinction between ordinary and complete stabilizer preservation, and the Choi characterization for channels, is developed in Ref.~\cite{Seddon2019}; below we give the convention bridge for the trace-decreasing maps used here.

The average success and conditional fidelity are
\begin{align}
 s_\delta(\cE)&=\int\Tr\cE(\rho_t(\psi)^{\otimes3})\,d\psi,
 \label{eq:success}\\
 F_\delta(\cE)&=
 \frac{\int\Tr[\psi\,\cE(\rho_t(\psi)^{\otimes3})]\,d\psi}
 {s_\delta(\cE)}.
 \label{eq:fidelity}
\end{align}
Equation~\eqref{eq:fidelity} weights the conditional output fidelity by the probability of success on each target. It is neither the unweighted average of individually normalized output fidelities nor a worst-case objective. The branch may depend on $\delta$ and the prescribed $s$, but not on the unknown $\psi$.

Use an unnormalized Choi operator with inputs $A=A_1A_2A_3$ preceding output $B$:
\begin{align}
 J&=\sum_{a,b=0}^7|a\rangle\langle b|\otimes\cE(|a\rangle\langle b|),
 \label{eq:choi}\\
 \cE(X)&=\Tr_A[(X^T\otimes I_2)J].
 \label{eq:choi-action}
\end{align}
Complete positivity and trace nonincrease give
\begin{equation}
 J\geq0,\qquad \Tr_B J\leq I_8.
 \label{eq:tni}
\end{equation}
All transposes refer to the fixed computational basis, so $X^T=X$, $Y^T=-Y$, and $Z^T=Z$ for Pauli matrices.

Condition~\eqref{eq:complete} implies
\begin{equation}
 J\in\Cs_4,\qquad J=\sum_\phi q_\phi\phi,\qquad q_\phi\geq0,
 \label{eq:cone}
\end{equation}
because applying $\cE$ to half of three normalized Bell pairs produces $J/8$. Conversely, tensor any stabilizer-cone input and its reference with $J$, and project the three corresponding pairs onto normalized Bell states. The contraction gives $(\id_k\otimes\cE)(X)/8$. Stabilizer tensor products, Bell postselection, and discarding preserve the positive cone. Positive rescaling proves~\eqref{eq:complete}. This establishes the branch-cone equivalence directly, without substituting a channel-robustness normalization for physical success.

Define the Haar operators with the input transpose already included:
\begin{align}
 \Qh(t)&=\int(\rho_t(\psi)^T)^{\otimes3}\otimes\psi\,d\psi,
 \label{eq:Q}\\
 \Rh(t)&=\int(\rho_t(\psi)^T)^{\otimes3}\otimes I_2\,d\psi.
 \label{eq:R}
\end{align}
Then $s_\delta(\cE)=\Tr(J\Rh)$ and the fidelity numerator is $\Tr(J\Qh)$. In particular, $\Tr J/8$ is the success on a maximally mixed joint input, which need not equal~\eqref{eq:success} for the correlated Haar ensemble.

\section{An exact three-copy bound}
\label{sec:proof}

\begin{theorem}
\label{thm:main}
For every $0<\delta<1$ and $0<s\leq1$, the maps in Sec.~\ref{sec:model} satisfy
\begin{equation}
 \max_{\cE:\,s_\delta(\cE)=s}F_\delta(\cE)=1-\frac{\delta}{2}.
 \label{eq:optimum}
\end{equation}
The maximum is attained by discarding two inputs and accepting the retained input independently with probability $s$.
\end{theorem}

We first establish the resource-dependent part of the witness. Operators on omitted tensor factors are identities, and $I$ without a subscript denotes $I_{16}$. Set
\begin{equation}
 \cB=\sum_{i=1}^3
 \bigl(X_{A_i}X_B-Y_{A_i}Y_B+Z_{A_i}Z_B\bigr).
 \label{eq:star}
\end{equation}

\begin{lemma}
\label{lem:star}
Every pure four-qubit stabilizer projector $\phi$ obeys $\Tr(\cB\phi)\leq3$. Consequently,
\begin{equation}
 \Tr[(\cB-3I)J]\leq0\quad\text{for all }J\in\Cs_4.
 \label{eq:starcone}
\end{equation}
\end{lemma}

\emph{Proof.}
A pure stabilizer projector is $\phi=2^{-4}\sum_{g\in G}g$, with $G$ a maximal commuting signed Pauli group. Pauli orthogonality gives expectations in $\{0,1,-1\}$; the Paulis with nonzero expectations commute pairwise.

Consider the nine unsigned operators $\sigma^a_{A_i}\sigma^a_B$. On different input edges, two of them commute precisely when their Pauli axes agree. On the same edge, different axes commute because both local factors anticommute. If a commuting subset contains two axes on one edge, it contains no term from another edge, and hence has size at most three. Otherwise it contains at most one term per edge, again at most three. Thus at most three summands in~\eqref{eq:star} have nonzero expectation, and every signed summand contributes at most one. The cone bound follows by nonnegative linear combination.\hfill$\square$

To obtain the full noise dependence, put
\begin{equation}
 H(t)=\Qh(t)-\frac{1+t}{2}\Rh(t),
 \label{eq:H}
\end{equation}
and introduce the input-pair operator
\begin{equation}
 \cA=\sum_{1\leq i<j\leq3}\sum_{a=x,y,z}
       \sigma^a_{A_i}\sigma^a_{A_j}.
 \label{eq:A}
\end{equation}
Odd Bloch moments vanish, and $\int r_ar_b\,d\psi=\delta_{ab}/3$. Expanding the three input factors gives the exact polynomials
\begin{align}
 \Rh(t)&=\frac18\left(I+\frac{t^2}{3}\cA\right),
 \label{eq:R-expand}\\
 \Qh(t)&=\frac1{16}\left(I+\frac{t}{3}\cB
                    +\frac{t^2}{3}\cA\right)+t^3\mathcal K,
 \label{eq:Q-expand}
\end{align}
where $\mathcal K$ is independent of $t$. Three inputs give degree at most three; these are not perturbative truncations. Paired input $Y$ transposes cancel in $\cA$, while a single input $Y$ produces the minus sign in~\eqref{eq:star}.

The constant and quadratic terms cancel in~\eqref{eq:H}. Its linear coefficient is
\begin{equation}
 L=\frac{\cB-3I}{48}.
 \label{eq:L}
\end{equation}
An odd cubic with linear coefficient $L$ is determined by $L$ and its value at one. It follows identically that
\begin{equation}
 H(t)=t(1-t^2)L+t^3H(1).
 \label{eq:interpolation}
\end{equation}
At the pure-input endpoint,
\begin{equation}
 H(1)=-\int(\psi^T)^{\otimes3}\otimes(I_2-\psi)\,d\psi\leq0.
 \label{eq:endpoint}
\end{equation}
Every integrand before the minus sign is a positive product operator. The argument does not use positivity of a partial transpose on an arbitrary entangled state.

For an explicit endpoint projector, the pure-qubit identity $I_2-\psi=Y\psi^T Y$ yields
\begin{align}
 H(1)&=-\frac15P,
 \label{eq:H1}\\
 P&=(I_8\otimes Y)\Pi_{\mathrm{sym},4}(I_8\otimes Y),
 \label{eq:P}
\end{align}
where $\Pi_{\mathrm{sym},4}$ projects onto the fully symmetric subspace of four qubits. Indeed,
\begin{equation}
 \int(\psi^T)^{\otimes4}\,d\psi=\frac15\Pi_{\mathrm{sym},4}.
 \label{eq:haar4}
\end{equation}
The integral is supported on the symmetric subspace, commutes with all fourfold qubit rotations, and has trace one. That subspace is the irreducible spin-two representation of dimension five, which fixes~\eqref{eq:haar4}. Thus $P^2=P\geq0$ and $\Tr P=5$. Equations~\eqref{eq:interpolation} and~\eqref{eq:H1} prove~\eqref{eq:certificate-intro}.

For $0\leq t\leq1$, both $t(1-t^2)$ and $t^3$ are nonnegative. By Lemma~\ref{lem:star}, $J\in\Cs_4$, and $J\geq0$,
\begin{equation}
 \Tr[JH(t)]\leq0.
 \label{eq:pairing}
\end{equation}
This pairing bound uses no particular value of success and no assumption that success is constant over target states. Dividing by the actual positive number $\Tr(J\Rh)=s$ proves
\begin{equation}
 F_\delta(\cE)\leq\frac{1+t}{2}=1-\frac{\delta}{2}.
 \label{eq:upper}
\end{equation}

For attainability, identify the first input with the output and take
\begin{equation}
 \cE_s(X)=s\Tr_{A_2A_3}(X).
 \label{eq:attainer}
\end{equation}
With tensor factors placed in input-first order, its Choi operator is
\begin{equation}
 J_s=2s\,\Phi^+_{A_1B}\otimes I_{A_2A_3},
 \qquad \Tr_BJ_s=sI_8,
 \label{eq:primal}
\end{equation}
where $\Phi^+$ is the normalized Bell projector. Expanding the two identity factors in the computational basis shows $J_s\in\Cs_4$. The operation also has a direct realization by discarding and independent classical acceptance. For every target, its output is $s\rho_t(\psi)$; hence its success is $s$ and its fidelity numerator is $s(1+t)/2$. This proves~\eqref{eq:optimum}.\hfill$\square$

\section{Exact verification and physical scope}
\label{sec:verification}

The preceding argument is analytic and does not require stabilizer enumeration. Exact finite-dimensional calculations supply separate checks on its signs, normalizations, and coverage. Two input reconstructions use permutation-symmetric Haar projectors and direct Bloch moments, respectively. They agree after conversion to the input-first Choi convention. The two-copy case $\delta=s=1/2$ recovers the known fidelity $3/4$ from Ref.~\cite{He2025NoGo}, Appendix~D; this is a baseline, not a new result.

An exhaustive construction extends binary isotropic Pauli subspaces and assigns all generator eigenvalues. It produces $2295$ maximal isotropic spaces and $36720$ pure four-qubit stabilizer projectors. A separate construction generates the Clifford orbit of $|0000\rangle$. The two constructions agree on every signed Pauli support and on every witness polynomial. An additional calculation using affine computational-basis supports and quadratic phases also agrees on all projectors and coefficients. Appendix~\ref{app:enumeration} explains the completeness and interval certificate.

All scientific arithmetic in these checks is exact rational or Gaussian-integer arithmetic. The largest dense operator is $16\times16$. In the integer projector calculation the intermediate real and imaginary magnitudes are bounded by $8192$, within signed 64-bit range; other exact calculations use arbitrary-precision integers. There is no noise grid, sampling error, numerical fitting, solver tolerance, or infinite-dimensional truncation. Code, exact input representations, polynomial data, and regeneration instructions accompany the draft. The stated test environment is Python~3.9.6, SymPy~1.14.0, and NumPy~2.0.2.

The operational implication is that postselection cannot improve the success-weighted universal mean fidelity using three such inputs within the class~\eqref{eq:complete}, however small the positive prescribed success rate. The input may contain nonstabilizerness; the limitation is on the successful operation. The first term of~\eqref{eq:certificate-intro} is bounded through stabilizer commutation, while the second is bounded by ordinary positivity. Consequently the resource restriction has a specific role in the proof.

The theorem is a statement about a finite-dimensional one-shot model with dimensionless noise strength, success, and fidelity. It concerns independent isotropic noise on three copies of the same unknown pure qubit. It supplies no claim for more copies, correlated or anisotropic noise, a different prior ensemble, or a worst-case fidelity. In particular, it does not settle the arbitrary-copy conjecture discussed in Ref.~\cite{He2025NoGo}. The simple cubic interpolation used here relies on having three noisy input factors.

No general free completion of every admissible success branch is assumed or proved. The attaining branch~\eqref{eq:attainer} can itself be paired with the complementary random rejection branch. Similarly, no assertion that all maps satisfying~\eqref{eq:complete} have Clifford-circuit realizations is needed. The continuous values at $t=0$ and $t=1$ give checks of the maximally noisy and noiseless limits, respectively; the theorem states the original open noise interval and never divides by zero success.

\subsection{Use of artificial intelligence}

OpenAI Codex was used for analytic derivation, generation of exact-computation code, literature comparison, and manuscript drafting. Separate AI sessions produced alternative input and stabilizer representations and checked the argument. These computations are described above and supplied as executable evidence; AI session separation is not external peer review or formal verification. Model-version identifiers were not recorded for all earlier sessions in the materials used for this draft. Human author review of the scientific argument, citations, code, and AI-assisted text, and final confirmation of the disclosure, remain pending.

\section*{Data availability}
The exact input operators, complete stabilizer-ray data, certificate summaries, and code are included in the accompanying reproducibility files. They are generated mathematical data, not experimental observations. The reproducibility guide explains how to regenerate outputs in a separate working directory. Third-party articles are cited by their public identifiers and are not included in that package.

\appendix
\section{Complete enumeration and interval checks}
\label{app:enumeration}

Unsigned $n$-qubit Paulis form $\mathbb F_2^{2n}$, with commutation given by the symplectic form. Starting from the zero space, extend each isotropic subspace $L$ by a vector in $L^\perp\setminus L$, deduplicating by the entire resulting space. Every isotropic space occurs: any $(k+1)$-dimensional isotropic space contains a $k$-dimensional one and a complementary vector. For $n=4$, the counts at dimensions zero through four are
\begin{equation}
 1,\quad255,\quad5355,\quad11475,\quad2295.
 \label{eq:counts}
\end{equation}
For each maximal space, choose four independent commuting Hermitian generators $P_j$. The $16$ joint eigenspace projectors are
\begin{equation}
 \phi_{\bm\epsilon}=\frac1{16}\prod_{j=1}^4(I+\epsilon_jP_j),
 \qquad\epsilon_j\in\{1,-1\}.
 \label{eq:projector-enum}
\end{equation}
The factors commute. Every nonempty generator product is a nonidentity Pauli, so each projector has trace one and rank one. Conversely, every pure stabilizer projector has such a maximal support and an eigenvalue assignment. Its support and eigenvalues are uniquely recovered from its nonzero Pauli expectations. This proves both completeness and lack of overcounting, rather than using the final count alone as the justification.

For the explicit ray polynomials, let $\eta_x=\eta_z=1$, $\eta_y=-1$, and set
\begin{align}
 \cC={}&\sum_{a,b,c,d}\eta_a\eta_b\eta_c
 (\delta_{ab}\delta_{cd}+\delta_{ac}\delta_{bd}
                         +\delta_{ad}\delta_{bc})\nonumber\\
 &\hspace{8mm}\times\sigma^a_{A_1}\sigma^b_{A_2}
                          \sigma^c_{A_3}\sigma^d_B,
 \label{eq:C}
\end{align}
where all four indices run over $x,y,z$. The fourth Bloch moment is the displayed delta sum divided by $15$, so $\mathcal K=\cC/240$. Equations~\eqref{eq:R-expand} and~\eqref{eq:Q-expand} give, for every ray,
\begin{align}
 g_\phi(t)&=\Tr[H(t)\phi]=\frac{t}{240}
                  (\alpha_\phi+\beta_\phi t^2),
 \label{eq:ray}\\
 \alpha_\phi&=5\Tr(\cB\phi)-15,\nonumber\\
 \beta_\phi&=\Tr(\cC\phi)-5\Tr(\cA\phi).
 \label{eq:coefficients}
\end{align}
The complete data contain $36$ coefficient pairs and $180$ identically zero ray polynomials. Every pair satisfies
\begin{equation}
 \alpha_\phi\leq0,\qquad\alpha_\phi+\beta_\phi\leq0.
 \label{eq:endpoints-ray}
\end{equation}
Since
\begin{equation}
 \alpha+\beta t^2=(1-t^2)\alpha+t^2(\alpha+\beta),
\end{equation}
these exact inequalities cover the entire interval $0\leq t\leq1$. The zero-polynomial count is a verification statistic, not a classification of all optimal physical protocols. The analytic proof in Sec.~\ref{sec:proof} establishes the operational bound independently of this enumeration.

\bibliography{references}
\end{document}
